\section{占有制度}
\begin{enumerate}
  \item 考非所有权人的占有，\zd{盗赃物和遗失物}，1 个选项。
  \item 保护对物的事实支配状态。
  \item 分类。
        \begin{enumerate}
          \item 有权、\zd{无权}。是否存在本权（权源），\zd{可相互转化}。
          \item 善意（\zd{必要费用请求权}）、恶意。\zd{都是无权占有}，区分标准不一样而已。\al{误拿他人行李}。
          \item 自主、他主。占有人的主观心态，以自己所有的意思而占有。想还 \yjt 他，不想 \yjt 自。
          \item 直接、间接。不看是不是在你手上，相对而言。\zd{占有媒介关系、他主占有、返还原物请求权}。
          \item 盗赃物和遗失物：无权、恶意、自主。
        \end{enumerate}
  \item 占有的保护。
        \begin{enumerate}
          \item 占有保护请求权，和物权请求权的区别，\zd{请求权基础不同}。
                \begin{enumerate}
                  \item \zd{占有返还请求权}，自侵占发生之日起 1 年。
                  \item 停止侵害请求权。
                  \item 排除妨碍请求权。
                  \item 消除危险请求权。
                \end{enumerate}
          \item 占有损害赔偿请求权。
        \end{enumerate}
  \item 占有的推定效力。\zd{占有即所有}，相对推定正确的效力，允许别人反证。
\end{enumerate}
